% GENERATED FILE. Edit canonical lessons, then run npm run generate:books.
% canonical-source-hash: fnv1a64:98507fe4c9ada79b
% canonical-lessons: ZH-C242-zhu, ZH-C242-shanyang, ZH-C242-she, ZH-C242-chongzi, ZH-C242-wenzi

\chapter{Farm Animals and Bugs}
\label{ch:zh-farm-animals-and-bugs}
\hlchaptermodality{242}

\begin{chapteropening}
\textbf{By the end of this chapter:} \emph{I can name a pig, a goat, a snake, insects and mosquitoes.}

Complete the last lesson of chapter 242: I can name a pig, a goat, a snake, insects and mosquitoes.
\end{chapteropening}

\section[zhū]{zhū — pig}

\label{lesson:ZH-C242-zhu}



\begin{quote}
Before the new one: say the Mandarin for to leave hospital, then the Mandarin for to recover.
\end{quote}

\subsection*{You'll want to know: zhū}
\emph{zhū} — \textquotedblleft{}pig\textquotedblright{}.

\begin{quote}
\emph{tā jiā yǒu sān zhī zhū} — His family keeps three pigs.
\end{quote}

Its characters come later; for now it is a word for the ear.

\begin{cousinweb}
\emph{zhū} is inherited from Middle Chinese \emph{trjo} (Guangyun, in Baxter's transcription), which glosses it, after an older dictionary, as the young of the pig.
\end{cousinweb}

\begin{grammarlens}[title={What you've built}]
One more word for everyday things.
\end{grammarlens}

\subsection*{Your turn}
\begin{itemize}
\raggedright
  \item \emph{Say it:} \emph{zhū}
  \item \emph{Say it:} \emph{zhū}, once more
  \item \emph{Say it:} \emph{tā jiā yǒu sān zhī zhū} — His family keeps three pigs.
  \item \emph{Recall:} say the Mandarin for to leave hospital, then the Mandarin for to recover, then say \emph{zhū} again
\end{itemize}

\begin{tcolorbox}[breakable,colback=teal!4,colframe=teal!35!black,title={Before you move on}]
What does zhū mean? (\textquotedblleft{}Pig\textquotedblright{}.) Say it once more.
\end{tcolorbox}
\section[shānyáng]{shānyáng — goat}

\label{lesson:ZH-C242-shanyang}



\begin{quote}
Before the new one: say the Mandarin for to recover, then the Mandarin for a pig.
\end{quote}

\subsection*{You'll want to know: shānyáng}
\emph{shānyáng} — \textquotedblleft{}goat\textquotedblright{}.

\begin{quote}
\emph{zhè zhī shānyáng hěn xǐhuan chī cǎo} — This goat loves eating grass.
\end{quote}

Its characters come later; for now it is a word for the ear.

\begin{cousinweb}
\emph{shānyáng} is built from \zh{山} 'hill' and \emph{yáng} 'sheep'. The first is inherited from Middle Chinese \emph{srean} (Guangyun, in Baxter's transcription); the dictionary glosses the second, \emph{yang} there, with the pair 'oxen and sheep'.
\end{cousinweb}

\begin{grammarlens}[title={What you've built}]
One more word for everyday things.
\end{grammarlens}

\subsection*{Your turn}
\begin{itemize}
\raggedright
  \item \emph{Say it:} \emph{shānyáng}
  \item \emph{Say it:} \emph{shānyáng}, once more
  \item \emph{Say it:} \emph{zhè zhī shānyáng hěn xǐhuan chī cǎo} — This goat loves eating grass.
  \item \emph{Recall:} say the Mandarin for to recover, then the Mandarin for a pig, then say \emph{shānyáng} again
\end{itemize}

\begin{tcolorbox}[breakable,colback=teal!4,colframe=teal!35!black,title={Before you move on}]
What does shānyáng mean? (\textquotedblleft{}Goat\textquotedblright{}.) Say it once more.
\end{tcolorbox}
\section[shé]{shé — snake}

\label{lesson:ZH-C242-she}



\begin{quote}
Before the new one: say the Mandarin for a pig, then the Mandarin for a goat.
\end{quote}

\subsection*{You'll want to know: shé}
\emph{shé} — \textquotedblleft{}snake\textquotedblright{}.

\begin{quote}
\emph{wǒ hěn pà shé} — I'm terrified of snakes.
\end{quote}

Its characters come later; for now it is a word for the ear.

\begin{cousinweb}
\emph{shé} is inherited from Middle Chinese \emph{zyae} (Guangyun, in Baxter's transcription), glossed 'a venomous creature'.
\end{cousinweb}

\begin{grammarlens}[title={What you've built}]
One more word for everyday things.
\end{grammarlens}

\subsection*{Your turn}
\begin{itemize}
\raggedright
  \item \emph{Say it:} \emph{shé}
  \item \emph{Say it:} \emph{shé}, once more
  \item \emph{Say it:} \emph{wǒ hěn pà shé} — I'm terrified of snakes.
  \item \emph{Recall:} say the Mandarin for a pig, then the Mandarin for a goat, then say \emph{shé} again
\end{itemize}

\begin{tcolorbox}[breakable,colback=teal!4,colframe=teal!35!black,title={Before you move on}]
What does shé mean? (\textquotedblleft{}Snake\textquotedblright{}.) Say it once more.
\end{tcolorbox}
\section[chóngzi]{chóngzi — insect; bug}

\label{lesson:ZH-C242-chongzi}



\begin{quote}
Before the new one: say the Mandarin for a goat, then the Mandarin for a snake.
\end{quote}

\subsection*{You'll want to know: chóngzi}
\emph{chóngzi} — \textquotedblleft{}insect; bug\textquotedblright{}.

\begin{quote}
\emph{wǒ nǚ'ér bú pà chóngzi} — My daughter isn't afraid of bugs.
\end{quote}

Its characters come later; for now it is a word for the ear.

\begin{cousinweb}
\emph{chóngzi} is built from \emph{chóng} 'insect' and the noun ending \zh{子}, both inherited from Middle Chinese, \emph{drjuwng} and \emph{tsiX} (Guangyun, in Baxter's transcription), which says, after an older dictionary, that creatures with feet are called by the first.
\end{cousinweb}

\begin{grammarlens}[title={What you've built}]
One more word for everyday things.
\end{grammarlens}

\subsection*{Your turn}
\begin{itemize}
\raggedright
  \item \emph{Say it:} \emph{chóngzi}
  \item \emph{Say it:} \emph{chóngzi}, once more
  \item \emph{Say it:} \emph{wǒ nǚ'ér bú pà chóngzi} — My daughter isn't afraid of bugs.
  \item \emph{Recall:} say the Mandarin for a goat, then the Mandarin for a snake, then say \emph{chóngzi} again
\end{itemize}

\begin{tcolorbox}[breakable,colback=teal!4,colframe=teal!35!black,title={Before you move on}]
What does chóngzi mean? (\textquotedblleft{}Insect; bug\textquotedblright{}.) Say it once more.
\end{tcolorbox}
\section[wénzi]{wénzi — mosquito}

\label{lesson:ZH-C242-wenzi}



\begin{quote}
Before the new one: say the Mandarin for a snake, then the Mandarin for a bug.
\end{quote}

\subsection*{You'll want to know: wénzi}
\emph{wénzi} — \textquotedblleft{}mosquito\textquotedblright{}.

\begin{quote}
\emph{xiàtiān zhèr yǒu wénzi} — There are mosquitoes here in summer.
\end{quote}

Its characters come later; for now it is a word for the ear.

\begin{cousinweb}
\emph{wénzi} is built from \emph{wén} 'mosquito' and the noun ending \zh{子}. The noun ending is inherited from Middle Chinese \emph{tsiX} (Guangyun, in Baxter's transcription); the dictionary lists the first, \emph{mjun} there, beside another form of the character.
\end{cousinweb}

\begin{grammarlens}[title={What you've built}]
One more word for everyday things.
\end{grammarlens}

\subsection*{Your turn}
\begin{itemize}
\raggedright
  \item \emph{Say it:} \emph{wénzi}
  \item \emph{Say it:} \emph{wénzi}, once more
  \item \emph{Say it:} \emph{xiàtiān zhèr yǒu wénzi} — There are mosquitoes here in summer.
  \item \emph{Recall:} say the Mandarin for a snake, then the Mandarin for a bug, then say \emph{wénzi} again
\end{itemize}

\begin{tcolorbox}[breakable,colback=teal!4,colframe=teal!35!black,title={Before you move on}]
What does wénzi mean? (\textquotedblleft{}Mosquito\textquotedblright{}.) Say it once more.
\end{tcolorbox}
